%% kapa: exam=1 round=33 year=2022 period=2 form=A subject=accounting subject_ko=회계학 seq=9 number=49
%% source: Q-Net 감정평가사 기출문제 (cst00309) artlSeq=5209997 file=raw/1008/5209997/2022년도 제33회 감정평가사 1차시험 2교시(감정평가 관계법규, 회계학).pdf
%% conversion: kordoc 4.24.51 -> tools/gen1.py (auto); review=proofread-sonnet; flags=none
\begin{question}{49}
재고자산 회계처리에 관한 설명으로 옳지 않은 것은?
\begin{choices}
  \item 생산에 투입하기 위해 보유하는 원재료 및 기타 소모품은 제품의 원가가 순실현가능가치를 초과할 것으로 예상되더라도 감액하지 아니한다.
  \item 생물자산에서 수확한 농림어업 수확물로 구성된 재고자산은 공정가치에서 처분부대원가를 뺀 금액으로 수확시점에 최초 인식한다.
  \item 재고자산을 현재의 장소에 현재의 상태로 이르게 하는 데 기여하지 않은 관리간접원가는 재고자산의 취득원가에 포함할 수 없다.
  \item 매입할인이나 매입금액에 대해 수령한 리베이트는 매입원가에서 차감한다.
  \item 개별법이 적용되지 않는 재고자산의 단위원가는 선입선출법이나 가중평균법을 사용하여 결정한다.
\end{choices}
\end{question}
\answer{1}
